\documentclass[11pt,a4paper]{article}

% --------------------------------------------------
% PACKAGES
% --------------------------------------------------
\usepackage[utf8]{inputenc}
\usepackage[T1]{fontenc}
\usepackage[french]{babel}
\usepackage{lmodern}
\usepackage{geometry}
\usepackage{amsmath,amssymb,amsthm}
\usepackage{array}
\usepackage{booktabs}
\usepackage{graphicx}
\usepackage{float}
\usepackage{hyperref}
\usepackage{xcolor}
\usepackage{enumitem}
\usepackage{listings}
\usepackage{fancyhdr}
\usepackage{titlesec}
\usepackage{caption}

\geometry{margin=2.5cm}
\setlength{\parskip}{0.5em}
\setlength{\parindent}{0pt}

% --------------------------------------------------
% HYPERREF
% --------------------------------------------------
\hypersetup{
    colorlinks=true,
    linkcolor=blue!60!black,
    urlcolor=blue!60!black,
    citecolor=blue!60!black,
    pdftitle={Projet Optimisation Discrète - VRPTW},
    pdfauthor={Nom1 - Nom2}
}

% --------------------------------------------------
% CODE STYLE
% --------------------------------------------------
\lstset{
    basicstyle=\ttfamily\small,
    breaklines=true,
    frame=single,
    columns=fullflexible,
    keepspaces=true
}

% --------------------------------------------------
% HEADER / FOOTER
% --------------------------------------------------
\pagestyle{fancy}
\fancyhf{}
\fancyhead[L]{Projet Optimisation Discrète}
\fancyhead[R]{VRPTW}
\fancyfoot[C]{\thepage}

% --------------------------------------------------
% TITLE FORMAT
% --------------------------------------------------
\titleformat{\section}{\large\bfseries}{\thesection.}{0.5em}{}
\titleformat{\subsection}{\normalsize\bfseries}{\thesubsection.}{0.5em}{}

% --------------------------------------------------
% DOCUMENT
% --------------------------------------------------
\begin{document}

\begin{titlepage}
    \centering
    {\scshape\LARGE Polytech Lyon \par}
    \vspace{1cm}
    {\scshape\Large Optimisation Discrète \par}
    \vspace{1.5cm}
    {\huge\bfseries Projet : Vehicle Routing Problem with Time Windows\par}
    \vspace{1.5cm}
    {\Large Rapport de projet\par}
    \vspace{2cm}
    
    \begin{tabular}{rl}
        Étudiants : & Nom Prénom \\
                    & Nom Prénom \\
        Enseignant : & Stéphane Bonnevay \\
        Date : & \today
    \end{tabular}
    
    \vfill
    
    \begin{minipage}{0.9\textwidth}
    \small
    Ce rapport présente notre travail sur la résolution du problème de tournées de véhicules avec fenêtres de temps (VRPTW) à l'aide de métaheuristiques. Le sujet demande de modéliser le problème, de structurer le code, de déterminer le nombre minimal de véhicules, de générer des solutions aléatoires, puis d'implémenter et comparer deux métaheuristiques.
    \end{minipage}
    
    \vfill
\end{titlepage}

\tableofcontents
\newpage

% --------------------------------------------------
\section{Introduction}
% --------------------------------------------------

Le projet étudié porte sur le \textbf{Vehicle Routing Problem with Time Windows (VRPTW)}. L'objectif est de déterminer un ensemble de tournées partant et revenant au dépôt, de manière à desservir chaque client exactement une fois, en respectant les contraintes du problème et en minimisant la distance totale parcourue.

Le sujet demande explicitement :
\begin{itemize}
    \item de modéliser le problème et de structurer le code ;
    \item de déterminer, pour chaque jeu de données, le nombre minimum de véhicules ;
    \item de créer un générateur aléatoire de solutions ;
    \item d'implémenter deux métaheuristiques ;
    \item de comparer les résultats obtenus ;
    \item de commencer sans tenir compte des fenêtres de temps, puis de les ajouter ensuite.
\end{itemize}

Dans un premier temps, nous considérons donc le problème \textbf{sans fenêtres temporelles}, ce qui revient à traiter un \textbf{Capacitated Vehicle Routing Problem (CVRP)}. Les fenêtres de temps seront prises en compte dans une seconde phase.

% --------------------------------------------------
\section{Organisation du projet}
% --------------------------------------------------

Afin de répondre progressivement aux différentes questions du sujet, nous avons découpé le travail en plusieurs phases.

\subsection{Phase 1 : Fondations du projet}
Cette phase regroupe :
\begin{itemize}
    \item la modélisation mathématique du problème ;
    \item la définition de la structure du code.
\end{itemize}

\subsection{Phase 2 : Détermination du nombre minimal de véhicules}
Cette phase est décomposée en deux sous-parties :
\begin{itemize}
    \item sans fenêtres de temps ;
    \item avec fenêtres de temps.
\end{itemize}

\subsection{Phase 3 : Génération de solutions initiales}
Le sujet demandant un générateur aléatoire, nous avons choisi de construire des \textbf{solutions aléatoires faisables}.

\subsection{Phase 4 : Première métaheuristique}
Cette partie sera complétée ultérieurement.

\subsection{Phase 5 : Deuxième métaheuristique}
Cette partie sera complétée ultérieurement.

\subsection{Phase 6 : Protocole expérimental et comparaison}
Cette partie sera complétée ultérieurement.

% --------------------------------------------------
\section{Modélisation du problème}
% --------------------------------------------------

\subsection{Description générale}
On considère un dépôt unique et un ensemble de clients. Chaque client possède :
\begin{itemize}
    \item des coordonnées cartésiennes ;
    \item une demande ;
    \item éventuellement une fenêtre de temps ;
    \item un temps de service.
\end{itemize}

Les véhicules sont supposés identiques et possèdent tous la même capacité maximale \(Q\).

L'objectif principal est de minimiser la distance totale parcourue. Le nombre de véhicules n'étant pas fixé dans le sujet, il doit être déterminé.

\subsection{Ensembles}
\[
V = \{0,1,\dots,n\}
\]
désigne l'ensemble des sommets, où :
\begin{itemize}
    \item \(0\) représente le dépôt ;
    \item \(C = \{1,\dots,n\}\) représente l'ensemble des clients.
\end{itemize}

\subsection{Paramètres}
\begin{itemize}
    \item \(d_{ij}\) : distance entre les sommets \(i\) et \(j\) ;
    \item \(q_i\) : demande du client \(i\) ;
    \item \(Q\) : capacité maximale d'un véhicule ;
    \item \(a_i\) : date de début de la fenêtre de temps du client \(i\) ;
    \item \(b_i\) : date de fin de la fenêtre de temps du client \(i\) ;
    \item \(s_i\) : durée de service chez le client \(i\).
\end{itemize}

\subsection{Variables de décision}
Une modélisation classique consiste à utiliser des variables binaires :
\[
x_{ij} =
\begin{cases}
1 & \text{si un véhicule va directement de } i \text{ vers } j, \\
0 & \text{sinon.}
\end{cases}
\]

Pour la version avec fenêtres de temps, on introduit également :
\[
t_i = \text{instant de début de service au client } i.
\]

\subsection{Fonction objectif}
L'objectif est de minimiser la distance totale parcourue :
\[
\min \sum_{i \in V} \sum_{j \in V,\, j \neq i} d_{ij} x_{ij}.
\]

\subsection{Contraintes principales}

\subsubsection{Visite unique de chaque client}
Chaque client doit être visité exactement une fois :
\[
\sum_{j \in V,\, j \neq i} x_{ij} = 1 \qquad \forall i \in C
\]
\[
\sum_{i \in V,\, i \neq j} x_{ij} = 1 \qquad \forall j \in C
\]

\subsubsection{Conservation du flot}
Si un véhicule arrive chez un client, il doit également en repartir :
\[
\sum_{j \in V} x_{ij} = \sum_{j \in V} x_{ji} \qquad \forall i \in C
\]

\subsubsection{Capacité}
Pour chaque véhicule, la somme des demandes des clients affectés à sa tournée ne doit pas dépasser \(Q\).

Cette contrainte sera traitée dans notre implémentation par le calcul explicite de la charge de chaque tournée.

\subsubsection{Fenêtres de temps}
Dans la seconde partie du projet, on ajoute :
\[
a_i \le t_i \le b_i
\]
et, pour tout déplacement de \(i\) vers \(j\),
\[
t_j \ge t_i + s_i + d_{ij}.
\]

\subsection{Remarque sur l'utilisation d'OR-Tools}
Nous avons choisi d'utiliser OR-Tools uniquement pour des \textbf{tests sur petites instances} et pour la validation de certaines hypothèses. En effet, sur des instances plus grandes, l'utilisation directe d'un solveur exact peut devenir trop coûteuse en temps de calcul. La résolution principale du projet repose donc sur des méthodes heuristiques et métaheuristiques.

% --------------------------------------------------
\section{Structure du code}
% --------------------------------------------------

La structure du code a été pensée de manière modulaire, afin de séparer clairement les rôles de chaque composant.

\begin{lstlisting}
src/
  main.py
  model.py
  io_utils.py
  evaluate.py
  solver.py
  neighbors.py
  meta/
    simulated_annealing.py
    tabu_search.py
\end{lstlisting}

\subsection{\texttt{main.py}}
Point d'entrée du programme. Il permet de charger une instance, de lancer l'algorithme choisi et d'afficher ou sauvegarder les résultats.

\subsection{\texttt{model.py}}
Ce fichier contient les structures de données principales :
\begin{itemize}
    \item client ;
    \item dépôt ;
    \item instance ;
    \item solution.
\end{itemize}

\subsection{\texttt{io\_utils.py}}
Ce module est chargé de la lecture des fichiers \texttt{.vrp} et de la sauvegarde des solutions au format JSON.

\subsection{\texttt{evaluate.py}}
Ce fichier regroupe les fonctions de calcul :
\begin{itemize}
    \item distance totale ;
    \item faisabilité d'une tournée ;
    \item respect de la capacité ;
    \item respect des fenêtres de temps.
\end{itemize}

\subsection{\texttt{solver.py}}
Interface de haut niveau permettant de lancer les différentes méthodes de résolution.

\subsection{\texttt{neighbors.py}}
Ce module contiendra les opérateurs de voisinage utilisés par les métaheuristiques, par exemple :
\begin{itemize}
    \item déplacement d'un client ;
    \item échange de clients ;
    \item inversion de segment (\texttt{2-opt}) ;
    \item opérateurs intra-route et inter-route.
\end{itemize}

\subsection{\texttt{meta/}}
Ce dossier contiendra les deux métaheuristiques retenues dans le projet :
\begin{itemize}
    \item recuit simulé ;
    \item recherche tabou.
\end{itemize}

% --------------------------------------------------
\section{Détermination du nombre minimal de véhicules}
% --------------------------------------------------

\subsection{Principe général}
Le sujet demande de déterminer, pour chaque jeu de données, le nombre minimum de véhicules à utiliser.

Nous avons choisi de traiter cette question en deux temps :
\begin{itemize}
    \item d'abord sans fenêtres de temps ;
    \item ensuite avec fenêtres de temps.
\end{itemize}

Dans cette première étape, nous ne considérons donc que :
\begin{itemize}
    \item la capacité des véhicules ;
    \item la nécessité de construire de vraies tournées partant et revenant au dépôt ;
    \item la distance entre les clients et le dépôt.
\end{itemize}

Autrement dit, nous traitons d'abord une version de type \textbf{CVRP} avant d'ajouter les fenêtres de temps.

\subsection{Borne inférieure théorique}
Une borne inférieure immédiate du nombre de véhicules est obtenue à partir de la somme des demandes :
\[
LB = \left\lceil \frac{\sum_{i \in C} q_i}{Q} \right\rceil.
\]

Cette borne est nécessaire mais pas toujours suffisante, car elle ne tient compte que de la capacité totale disponible. Elle ignore :
\begin{itemize}
    \item la répartition des demandes entre les tournées ;
    \item la structure géographique des clients ;
    \item la cohérence réelle des trajets.
\end{itemize}

Deux clients compatibles en termes de charge peuvent en effet être très éloignés, ce qui rend leur regroupement peu pertinent dans une même tournée.

\subsection{Choix méthodologique}
Nous avons décidé de ne pas déterminer le nombre minimal de véhicules uniquement à partir des demandes, comme dans un simple problème de remplissage. En effet, même sans fenêtres temporelles, il est nécessaire de tenir compte de la distance afin de construire des tournées réalistes.

Notre méthode consiste donc à :
\begin{enumerate}[label=\arabic*.]
    \item calculer la borne inférieure \(LB\) ;
    \item fixer \(k = LB\) ;
    \item essayer de construire une solution faisable avec exactement \(k\) véhicules ;
    \item si aucun ensemble de tournées faisables n'est obtenu, incrémenter \(k\) ;
    \item retenir le premier \(k\) pour lequel tous les clients sont servis.
\end{enumerate}

\subsection{Heuristique utilisée pour tester une valeur de \(k\)}
Pour tester si \(k\) véhicules suffisent, nous utilisons une heuristique gloutonne géographique sous contrainte de capacité.

Le principe est le suivant :
\begin{enumerate}[label=\arabic*.]
    \item initialiser \(k\) tournées vides ;
    \item choisir un premier client pour chaque tournée ;
    \item compléter chaque tournée en ajoutant à chaque étape un client non encore servi ;
    \item parmi les clients admissibles, sélectionner celui qui minimise la distance au dernier client de la tournée ;
    \item ne jamais dépasser la capacité maximale \(Q\) ;
    \item s'arrêter lorsque plus aucun client ne peut être ajouté à la tournée courante.
\end{enumerate}

Un client est dit \textbf{admissible} si son ajout à la tournée ne fait pas dépasser la capacité du véhicule.

\subsection{Critère de faisabilité}
Pour une valeur donnée de \(k\), la solution est considérée comme faisable si :
\begin{itemize}
    \item tous les clients sont affectés à une unique tournée ;
    \item la charge de chaque tournée est inférieure ou égale à \(Q\).
\end{itemize}

Dans cette phase, la distance intervient comme \textbf{critère de construction} des tournées, mais pas encore comme objectif d'optimisation final.

\subsection{Remarque}
Cette approche ne prouve pas mathématiquement que le nombre trouvé est optimal. En revanche, elle fournit une estimation cohérente du nombre minimal de véhicules, adaptée à la taille des instances étudiées. Cette estimation sera ensuite affinée indirectement par les métaheuristiques développées dans les phases suivantes.

\subsection{Résultats obtenus sans fenêtres de temps}
Nous avons appliqué notre méthode de construction gloutonne géographique à l'ensemble des instances, en commençant à chaque fois par la borne inférieure
\[
LB = \left\lceil \frac{\sum_{i \in C} q_i}{Q} \right\rceil.
\]

Pour chaque instance testée, une solution faisable a été obtenue dès cette borne inférieure. Cela signifie que, dans le cas sans fenêtres de temps, le nombre minimal de véhicules est exactement égal à \(LB\) pour toutes les instances considérées.

\begin{center}
\begin{tabular}{lcccc}
\toprule
Instance & Capacité & Demande totale & \(LB\) & Nombre minimal trouvé \\
\midrule
data101.vrp  & 200  & 1458 & 8 & 8 \\
data102.vrp  & 200  & 1458 & 8 & 8 \\
data1101.vrp & 200  & 1724 & 9 & 9 \\
data1102.vrp & 200  & 1724 & 9 & 9 \\
data111.vrp  & 200  & 1458 & 8 & 8 \\
data112.vrp  & 200  & 1458 & 8 & 8 \\
data1201.vrp & 1000 & 1724 & 2 & 2 \\
data1202.vrp & 1000 & 1724 & 2 & 2 \\
data201.vrp  & 1000 & 1458 & 2 & 2 \\
data202.vrp  & 1000 & 1458 & 2 & 2 \\
\bottomrule
\end{tabular}
\end{center}

On remarque donc que, pour ces jeux de données et sans tenir compte des fenêtres temporelles, la contrainte de capacité suffit à expliquer le nombre minimal de véhicules.

\subsection{Influence des fenêtres de temps sur la borne inférieure}

La borne inférieure utilisée dans notre étude est basée uniquement sur la capacité des véhicules :

\[
LB = \left\lceil \frac{\sum_{i \in C} q_i}{Q} \right\rceil.
\]

Cette borne ne dépend que de la demande totale et de la capacité des véhicules. Elle est donc indépendante des distances et des fenêtres temporelles.

Ainsi, l'introduction des fenêtres de temps ne modifie pas cette borne inférieure. En revanche, elle restreint fortement l'ensemble des tournées faisables. Certaines combinaisons de clients compatibles en termes de capacité peuvent devenir impossibles à réaliser lorsque leurs fenêtres de service sont incompatibles.

Par conséquent, le nombre minimal de véhicules nécessaire avec fenêtres de temps est généralement supérieur ou égal à la borne inférieure :

\[
k_{TW} \ge LB.
\]

L'objectif de la seconde partie de la question 2 est donc de déterminer expérimentalement le nombre minimal de véhicules lorsque les contraintes temporelles sont prises en compte.

% --------------------------------------------------
\section{Génération d'une solution initiale}
% --------------------------------------------------

Le sujet demande explicitement la création d'un générateur aléatoire de solutions. Nous avons donc choisi de construire une \textbf{solution aléatoire faisable}.

\subsection{Choix retenu}
Nous ne retenons pas une permutation totalement aléatoire des clients, car celle-ci produit très souvent des solutions infaisables. À la place, nous construisons la solution progressivement :

\begin{enumerate}[label=\arabic*.]
    \item ouvrir une nouvelle tournée ;
    \item déterminer la liste des clients admissibles dans cette tournée ;
    \item choisir aléatoirement l'un de ces clients ;
    \item l'ajouter à la tournée ;
    \item recommencer jusqu'à ce qu'aucun autre client admissible ne puisse être ajouté ;
    \item ouvrir une nouvelle tournée ;
    \item répéter jusqu'à ce que tous les clients soient servis.
\end{enumerate}

\subsection{Cas sans fenêtres de temps}
Dans un premier temps, un client est admissible si son ajout ne fait pas dépasser la capacité du véhicule.

\subsection{Cas avec fenêtres de temps}
Dans un second temps, un client sera admissible si son ajout respecte à la fois :
\begin{itemize}
    \item la capacité du véhicule ;
    \item la faisabilité temporelle de la tournée.
\end{itemize}

\subsection{Rôle de cette solution}
Cette solution initiale n'a pas vocation à être optimale. Elle sert de point de départ aux métaheuristiques.

% --------------------------------------------------
\section{Choix des méthodes de résolution}
% --------------------------------------------------

Le projet demandant l'utilisation de deux métaheuristiques, nous avons retenu les approches suivantes :
\begin{itemize}
    \item le recuit simulé ;
    \item la recherche tabou.
\end{itemize}

Ces deux méthodes sont particulièrement adaptées à une résolution par voisinage du problème de tournées de véhicules.

\subsection{Voisinages envisagés}
Les opérateurs suivants seront étudiés :
\begin{itemize}
    \item relocate ;
    \item exchange ;
    \item 2-opt ;
    \item éventuellement des opérateurs inter-route plus complexes.
\end{itemize}

% --------------------------------------------------
\section{Format des données}
% --------------------------------------------------

Les fichiers d'entrée sont des fichiers texte au format \texttt{.vrp}. Une instance contient :
\begin{itemize}
    \item le nom du fichier ;
    \item le nombre de dépôts ;
    \item le nombre de clients ;
    \item la capacité maximale d'un véhicule ;
    \item les informations sur le dépôt ;
    \item les informations sur chaque client.
\end{itemize}

Exemple de structure :
\begin{lstlisting}
NAME: data101.vrp
COMMENT:
TYPE: vrptw
COORDINATES: cartesian
NB_DEPOTS: 1
NB_CLIENTS: 100
MAX_QUANTITY: 200

DATA_DEPOTS [idName x y readyTime dueTime]:
d1 35 35 0 230

DATA_CLIENTS [idName x y readyTime dueTime demand service]:
c1 41 49 161 171 10 10
...
\end{lstlisting}

Chaque client est donc décrit par :
\begin{itemize}
    \item ses coordonnées \((x,y)\) ;
    \item sa fenêtre de temps \([readyTime, dueTime]\) ;
    \item sa demande ;
    \item sa durée de service.
\end{itemize}

% --------------------------------------------------
\section{Travail restant}
% --------------------------------------------------

Les parties suivantes restent à compléter :
\begin{itemize}
    \item implémentation détaillée du calcul de faisabilité ;
    \item calcul expérimental du nombre minimal de véhicules sur chaque instance ;
    \item définition précise des voisinages ;
    \item implémentation du recuit simulé ;
    \item implémentation de la recherche tabou ;
    \item protocole expérimental ;
    \item analyse comparative des résultats ;
    \item prise en compte complète des fenêtres de temps.
\end{itemize}

% --------------------------------------------------
\section{Conclusion provisoire}
% --------------------------------------------------

À ce stade, nous avons :
\begin{itemize}
    \item identifié clairement le problème à résoudre ;
    \item proposé une première modélisation mathématique ;
    \item structuré les grandes phases du projet ;
    \item défini l'architecture logicielle ;
    \item choisi d'utiliser OR-Tools uniquement sur petites instances ;
    \item retenu une stratégie de génération aléatoire faisable pour construire les solutions initiales.
\end{itemize}

La suite du travail consistera à déterminer le nombre minimal de véhicules, puis à développer et comparer deux métaheuristiques adaptées au problème.

\end{document}